\section{証明の範囲と残る境界の問題}

正共鳴の全整数にわたる結果は、指数半群、Ore表示、正規スライス、有限整拡大による。
原点根基の冪零指数には全重みでの所属恒等式と非所属評価を与えた。
量子・Poisson素スペクトルの一致は、正規元の内部自己同型から局所冪零導分を回収する補間による。
葉の閉包公式は全域で成り立つが、無限次元単純加群そのものの全同型類や部分解消上の代数層は別の問題である。

一の冪根では、巡回ノルム、有限エタール被覆上の重み冪等元、非零辺の縮約により全位数を扱った。
行列代数になる開部分では、閉点ファイバーの単純性に加えて局所自由性も示す。
非Azumayaファイバーの全体は、道の簡約による上界と明示的な行列表現による下界を合わせて決定する。
特異部分への二次降下と一般の巡回降下も、互いに逆な行列表現を構成して証明した。
計算コードは有限個の位数についてこれらの全関係を厳密に照合するが、
全位数への結論には本文の一般証明を用いる。

Brauer周期と指数は、整域性と巡回ノルムの付値によって決まり、有限位数の標本からの推測ではない。
分裂被覆上で得られるのはベクトル束の自己準同型環である。
元の中心上では非零のBrauer類が残るため、通常の自己準同型環としてのNCCRと
同一視できない。この障害は、既知の中心開部分の閉点での完成局所McKay模型と両立する。
さらに全中心の最深点では、完成後にも指数 $\ell$ が残ることを独立した付値計算で証明した。

$U\ne0$ の中心は混合項 $H$ の剰余 $h$ をノルムとして記憶する。
この復元の結論は $Q,U,J$ を固定したものである。
中心元を $H$ に加える操作はOre導分を変えず、$J$ の移動に対応する。
全中心を求めた $f=0$ と $f=Qu^3$ の例では、特異部分の次元から
標識を忘れた全中心の非同型性も従う。
ここからアフィン代数のMorita・dg導来Morita非同値を得るが、射影導来圏の非同値までは結論しない。

混合項が $f(Q,u)+xG$ となる族については全中心を決定した。
混合項が中心元と $xG$ の和になる部分族では全閉点ファイバーも分類し、
$f\in Q\mathbb C[Q,u]$ では最深点のファイバーと完備局所Brauer障害が同じであることを示した。
一般境界の全形式模型からは、単純加群の巡回被覆と例外因子のBrauer留数を対応させた。
ここで用いるMcKay/NCCRの一般理論と巡回代数の留数公式は既存の理論であり、
本稿の内容はこの量子化族における明示表示と、その大域・完備局所的な相違である。

残る対象は、この一変数形に帰着しない混合項 $H$ に対する全中心と深い境界、
追加衝突点の全形式代数、および標識を忘れた射影圏の同値である。
最深点の代数自身に対するEnd型の障害は、中心が別のNCCRを持たないという結論ではない。
中心局所化や係数体上のBrauer障害を、任意の導来埋込みの不存在へ読み替えない。
